% SOURCES: docs/0.2.0v/plan.md (§4, §6, §17, §7); CLAUDE.md (Sprint Model, Git Workflow);
%          .claude/rules/*; docs/1.0.0v/supervisor-report.md (§8)

\chapter{Work Methodology}
\label{ch:methodology}

This chapter describes \emph{how} the team worked --- the development model,
ownership structure, interface contracts, and quality gates --- rather than what
was built. The process was unusually well documented for a student project, and
that discipline is itself one of the project's results.

\section{Development Model: Task-Driven Milestones}
Development followed a \emph{task-driven, milestone-gated} model rather than a
fixed time-box. Progress was unlocked by completing the tasks that gate each
milestone, not by the calendar: a milestone was ``done'' only when its gating tasks
were complete \emph{and} every quality gate was green. Five milestones carried the
work from \textbf{M1 (Foundation)} through \textbf{M5 (Polish + Demo)}, followed by
the \textbf{v1.0.0} demo cut (the milestone table is reproduced in
Chapter~\ref{ch:plan}, Table~\ref{tab:milestones}). Every task carried a single
owner, a priority (\texttt{P0} blocks the demo, \texttt{P1} is a Tier-1 feature,
\texttt{P2} is polish), and explicit \texttt{depends:}/\texttt{blocks:} links, so
the critical path was always visible and the Document Model --- which blocks every
other area --- was scheduled to land first.

\section{Team Structure and Ownership}
% SOURCES: CLAUDE.md (Module Ownership); docs/0.2.0v/plan.md (§10 Tasks)
Three engineers owned disjoint areas of the codebase, which kept merge contention
low and made every file traceable to a responsible owner:

\begin{itemize}[leftmargin=1.5em]
  \item \textbf{Ibrahim Haski} --- the Document Model
        (\texttt{types}, \texttt{schemas}, \texttt{operations}, \texttt{tokens},
        \texttt{validation}, \texttt{migrations}, \texttt{presets}), the output
        pipeline (\texttt{generator}, \texttt{runtime}, \texttt{seo},
        \texttt{export}), the Electron shell and IPC, and the build/CI
        configuration.
  \item \textbf{Yousef Deeb} --- the state layer: the Zustand stores, undo/redo
        history implemented as \texttt{immer} patches, persistence
        (the \texttt{.dtw} file format), autosave and crash recovery, and
        external-change reload; and, in v1.0.0, the draw / match / MCP surfaces.
  \item \textbf{Abd Alrazak Qahwaji} --- the user interface: the canvas, the
        properties, layers, and tokens panels, the sidebar, the topbar, and the
        semantic-inference hints the generator consumes.
\end{itemize}

\section{Cross-Owner Contracts}
% SOURCES: docs/0.2.0v/plan.md (§6 Contracts C1–C12); docs/0.2.0v/architecture.md (§6)
Because three people worked in parallel against a shared model, the interfaces
between owners were fixed as twelve numbered \emph{contracts} (C1--C12). These
include the shared types (C1), the Zod schemas kept in lockstep with them (C2), the
operation set (C3), the typed IPC surface (C4), the store hooks (C5), the generator
entry point (C6), the presets registry (C7), the validation function (C8), the
token resolver (C9), the semantic-inference hint (C10), the image-upload IPC
contract (C11), and the export orchestrator (C12). Any breaking change to a
contract required a pull request labelled \texttt{contract-change} and a review
from the named downstream consumer; this single rule is what allowed the three
lanes to move quickly without silently diverging.

\section{Quality Gates and Tooling}
% SOURCES: .claude/rules/git.md, preflight.md; CLAUDE.md (Commands, Code Quality)
Every push had to pass three gates: linting (ESLint + Prettier), strict
type-checking (TypeScript in strict mode across the main, preload, and renderer
\texttt{tsconfig}s), and the full automated test suite. These were enforced twice
--- locally by a \texttt{husky} pre-push hook and again in continuous integration
--- so a red build could not reach \texttt{main}. On top of the three gates,
specialised automated review passes checked properties that unit tests do not
naturally cover: output \emph{determinism}, the no-\texttt{position:absolute}
invariant, accessibility of generated output, and correctness of any contract
change. A weekly integration ritual merged all open feature branches, ran the full
matrix, and held a short debrief with no new commits, which caught cross-lane
integration problems early rather than at release. As a matter of academic
integrity, the project also enforced a ``no AI attribution'' policy on all commits,
pull requests, and reviews; the work and its authorship are the students' own.

\section{Version Control and Release Workflow}
% SOURCES: .claude/rules/git.md; .github/workflows/{ci,release}.yml; CHANGELOG.md
Work used \emph{conventional commits} (for example
\texttt{feat(generator): emit skip-to-content link}) on short-lived feature
branches that were squash-merged into \texttt{main}; the \texttt{main} branch was
never force-pushed. Releases were tag-driven: pushing a \texttt{v*} tag triggered a
per-operating-system build matrix that produced native installers (Windows NSIS,
Linux AppImage/deb, macOS dmg/zip) and published a single GitHub Release. Versions
progressed v0.1.0 $\rightarrow$ v0.2.0 $\rightarrow$ v0.3.0 $\rightarrow$
\textbf{v1.0.0}, with each tag corresponding to a completed milestone.
